\section{ماژول modegpt\_compressor (MoDeGPT Compressor)}
\label{sec:module_modegpt_compressor}

\subsection{معرفی و انگیزه}
این ماژول در چارچوب پژوهش حاضر پیاده‌سازی و ارزیابی شده و فاقد مقاله مستقل جداگانه در ادبیات تحقیق است.

در بلوک‌های استاندارد ترانسفورمر، لایه‌های خطی متعددی وجود دارند که یا دارای بعد ورودی مشترک هستند (مانند پروجکشن‌های پرسش و کلید $W_q$ و $W_k$ در مکانیزم توجه، یا نگاشت‌های ورودی و دروازه $W_{\text{gate}}$ و $W_{\text{up}}$ در لایه‌های \lr{SwiGLU/MLP}) یا اطلاعات بافت زبانی را به صورت موازی پردازش می‌کنند. اعمال مستقل تجزیه مقادیر منفرد (\lr{SVD}) بر روی تک‌تک این ماتریس‌ها باعث نادیده گرفتن همبستگی و همپوشانی زیرفضاهای ویژگی میان آن‌ها شده و افزونگی پارامتری بالایی را به همراه دارد.

الگوریتم \lr{MoDeGPT (Modular Decomposition of Transformer Blocks)} این چالش را با تجزیه توأم (\lr{Joint Decomposition}) ماتریس‌های مکمل حل می‌کند. در این روش، ماتریس‌های پروجکشن مرتبط در امتداد بعد خروجی به یکدیگر پیوند داده شده و از طریق یک پایه متعامد کم‌مرتبه مشترک فاکتورگیری می‌شوند؛ در نتیجه نه‌تنها پارامترها به صورت بهینه کاهش می‌یابند، بلکه تقارن فضایی بازنمایی‌ها حفظ می‌گردد.

\subsection{مبانی نظری و فرمول‌بندی ریاضی}
فرض کنید دو ماتریس وزن خطی متعلق به یک بلوک ترانسفورمر که دارای بعد ورودی مشترک $d_{\text{in}}$ هستند با $\mathbf{W}_1 \in \mathbb{R}^{d_1 \times d_{\text{in}}}$ و $\mathbf{W}_2 \in \mathbb{R}^{d_2 \times d_{\text{in}}}$ مشخص شوند. در الگوریتم \lr{MoDeGPT}، این دو ماتریس در امتداد بعد خروجی ادغام می‌شوند:
\begin{equation}
\mathbf{W}_{\text{joint}} = \begin{bmatrix} \mathbf{W}_1 \\ \mathbf{W}_2 \end{bmatrix} \in \mathbb{R}^{(d_1 + d_2) \times d_{\text{in}}}
\end{equation}
سپس تجزیه مقادیر منفرد فشرده با رتبه بهینه $r \le \min(d_1 + d_2, d_{\text{in}})$ بر روی ماتریس مشترک اعمال می‌گردد:
\begin{equation}
\mathbf{W}_{\text{joint}} \approx \mathbf{U}_r \mathbf{\Sigma}_r \mathbf{V}_r^\top = \tilde{\mathbf{U}}_r \mathbf{V}_r^\top
\end{equation}
که در آن $\tilde{\mathbf{U}}_r = \mathbf{U}_r \mathbf{\Sigma}_r \in \mathbb{R}^{(d_1 + d_2) \times r}$ و $\mathbf{V}_r^\top \in \mathbb{R}^{r \times d_{\text{in}}}$ است. زیرفضای راست $\mathbf{V}_r^\top$ به عنوان یک پایه مشترک عمل می‌کند که ورودی را به یک فضای نهفته کم‌بعد تصویر می‌نماید. ماتریس چپ به دو بلوک متناظر با ابعاد اولیه تفکیک می‌شود:
\begin{equation}
\tilde{\mathbf{U}}_r = \begin{bmatrix} \tilde{\mathbf{U}}_1 \\ \tilde{\mathbf{U}}_2 \end{bmatrix}, \quad \tilde{\mathbf{U}}_1 \in \mathbb{R}^{d_1 \times r}, \quad \tilde{\mathbf{U}}_2 \in \mathbb{R}^{d_2 \times r}
\end{equation}
در نتیجه تقریب بازسازی‌شده هر یک از لایه‌ها به فرم زیر به دست می‌آید:
\begin{equation}
\hat{\mathbf{W}}_1 = \tilde{\mathbf{U}}_1 \mathbf{V}_r^\top, \quad \hat{\mathbf{W}}_2 = \tilde{\mathbf{U}}_2 \mathbf{V}_r^\top
\end{equation}
برای هر لایه، خطای نسبی فروبنیوس ارزیابی شده و با آستانه بحرانی $\tau_{\text{error}}$ مقایسه می‌گردد:
\begin{equation}
\mathcal{E}_i = \frac{\|\mathbf{W}_i - \hat{\mathbf{W}}_i\|_F}{\|\mathbf{W}_i\|_F + \epsilon}
\end{equation}

\subsection{معماری و پیاده‌سازی}
موتور فشرده‌سازی \lr{MoDeGPT} در فایل \lr{Codes/framework/compression\_engine/modegpt\_compressor.py} تعبیه شده و شامل مؤلفه‌های زیر است:
\begin{itemize}
    \item \textbf{لایه خطی \lr{MoDeGPTLoRALinear}:} لایه جایگزین که وزن‌های پایه فاکتورگیری‌شده را در فرمت استاندارد نگهداری کرده و با لایه‌های \lr{Conv1D} در مدل‌های سبکی همچون \lr{GPT-2} سازگاری کامل دارد. همچنین در صورت نیاز به فاین‌تیونینگ، قابلیت اضافه کردن ادپترهای \lr{LoRA} فعال می‌گردد.
    \item \textbf{کلاس \lr{MoDeGPTCompressor}:} مدیریت خوشه‌بندی لایه‌ها به جفت‌های مکمل (\lr{Joint Pairs})، بررسی سازگاری ابعاد لایه‌ها، مدیریت لایه‌های منفرد که جفت مناسبی ندارند توسط \lr{\_decompose\_single}، محافظت از پروجکشن‌های حساس در معماری‌های با مکانیزم توجه نامتقارن (\lr{GQA-aware skipping}) و امکان بازگشت امن به وزن اصلی در صورت نقض آستانه خطا.
\end{itemize}

\subsection{نتایج تجربی}
عملکرد تجربی فشرده‌ساز \lr{modegpt\_compressor} بر اساس گزارش بنچ‌مارک رسمی در \lr{results/unified\_modegpt\_4models\_summary.json} در جدول \ref{tab:benchmark_modegpt_compressor} خلاصه شده است:

\begin{table}[H]
\centering
\small
\caption{نتایج تجربی و اعتبارسنجی بنچ‌مارک فشرده‌ساز \lr{modegpt\_compressor} روی مدل‌های پایه}
\label{tab:benchmark_modegpt_compressor}
\begin{tabular}{lccccc}
\toprule
\textbf{مدل پایه} \lr{(Base Model)} & \textbf{دقت} \lr{(Accuracy)} & \textbf{شاخص} \lr{(F1-Score)} & \textbf{حافظه} \lr{(Peak VRAM)} & \textbf{زمان} \lr{(Time)} & \textbf{وضعیت} \\
\midrule
\lr{DistilBERT} & \lr{90.48\%} & \lr{0.9077} & \lr{311.2 MB} & \lr{17.9 s} & \lr{PASS} \\
\lr{RoBERTa} & \lr{72.02\%} & \lr{0.7724} & \lr{301.7 MB} & \lr{25.5 s} & \lr{PASS} \\
\lr{GPT-2} & \lr{60.55\%} & \lr{0.6861} & \lr{639.7 MB} & \lr{25.6 s} & \lr{PASS} \\
\lr{Qwen} & \lr{---} & \lr{---} & \lr{---} & \lr{---} & \lr{SKIP} \\
\bottomrule
\end{tabular}
\end{table}

بررسی داده‌های تجربی نشان می‌دهد که ماژول \lr{modegpt\_compressor} در مدل انکودری \lr{DistilBERT} دقت قابل توجه پایدار \lr{90.48\%} و شاخص F1 معادل \lr{0.9077} را بدون کوچکترین افت ثبت کرده است. در مدل \lr{RoBERTa} دقت \lr{72.02\%} حاصل شده و در مدل دیکودری \lr{GPT-2} دقت مدل فشرده‌شده به \lr{60.55\%} رسیده که افتی ناچیز در حد \lr{-0.23\%} نسبت به مدل پایه دارد و با موفقیت کامل وضعیت \lr{PASS} را احراز کرده است. در مدل \lr{Qwen} اجرای آزمون به دلیل ناهمگونی ساختار جفت‌های خطی چندسر به صورت \lr{SKIP} گزارش شده است.

\subsection{محدودیت‌ها و تحلیل شکست}
تحلیل فنی فرآیند فشرده‌سازی با \lr{modegpt\_compressor} نشان‌دهنده نکات و چالش‌های ساختاری زیر است:
\begin{itemize}
    \item \textbf{پیچیدگی جفت‌سازی در معماری‌های ناهمگن:} پیش‌فرض الگوریتم \lr{MoDeGPT} وجود جفت‌های متقارن خطی با ابعاد ورودی یکسان است. در مدل‌های مدرنی که از \lr{GQA} استفاده می‌کنند، ابعاد پروجکشن کلید و مقدار بسیار کوچکتر از پروجکشن پرسش است که امکان الحاق مستقیم بدون پدینگ یا نگاشت‌های اضافی را از بین برده و منجر به انصراف موتور از فشرده‌سازی جفت می‌گردد.
    \item \textbf{تثبیت رتبه سراسری بدون آگاهی از حساسیت طیفی:} تنظیم یک رتبه ثابت (\lr{rank=64}) برای کلیه زوج لایه‌ها بدون در نظر گرفتن محتوای اطلاعاتی طیف مقادیر منفرد می‌تواند در لایه‌های میانی که دربرگیرنده بیشترین پیچیدگی معنایی هستند منجر به افت دقت گردد.
\end{itemize}
